% v19: current support consequences. Six necessary types; no existence claim.
\subsection{葉の全排除と，残る六型}
\label{sf19:section}

\begin{corollary}[非空支持は二点連結で最小次数二以上]
\label{sf19:no-leaves}
正の局所整数 charts，marked puncture の整数 clutching，対向 $D$ の相殺，
および node 非負性を満たす，五次の標準 $K_5$ branch graph を考える．
高分母 support $H=\{ij:q_{ij}>1\}$ が非空なら，
$H$ は五頂点全てを使う二点連結 graph であり，最小次数は二以上である．
さらに
\[
 |E(H)|\ge7,\qquad b_1(H)\ge3
\]
であり，各二価頂点は三角形に属する．
特に，一つの core の高分母 ray が高々一本なら，全ての分母は一である．
\end{corollary}

\begin{proof}
系~\ref{kg17:low-core-propagation} により，非空の場合に孤立頂点はない．
木成分，単純閉路成分，二価二頂点を持つ三角形，葉付き四角形，
$K_{2,3}$，diamond，二葉三角形および四使用頂点 $K_4$ の障害を適用すると，
五頂点上の十一型・三百五十六標識が残る．
定理~\ref{hs18:no-house}，\ref{gten18:no-support}，\ref{lh18:no-K4-leaf} により，
そこから $G_{5,5},G_{5,10},G_{5,12}$ が除かれる．
八つの残存型で葉を持つのは $G_{5,7},G_{5,9}$ だけであり，
六辺の型もこの二つだけである．
定理~\ref{g719:no-G5-7} と定理~\ref{dl19:no-G5-9} がそれらを排除するので，
最小次数二以上，辺数七以上，$b_1(H)=|E(H)|-4\ge3$ を得る．
全ての残存型の二価頂点の二近傍は高辺で結ばれる．
これらの有限 graph の確認は次の表と全標識の検算に明記する．

二点連結性も初等的に従う．切断頂点 $v$ があるなら，
$H-v$ の各成分は最小次数二の条件により二頂点以上を持つ．
残る頂点は四つなので，二頂点ずつの二成分だけとなる．
各頂点の次数を二以上にするため，各成分は一本の辺を持ち，
その両端は $v$ に接続しなければならない．
これは二つの三角形が $v$ だけを共有する六辺の graph となり，辺数七以上に反する．
最後の結論は，孤立頂点と葉のいずれも許されないことの言い換えである．
\end{proof}

\begin{corollary}[残る六型と，低辺の補 graph]
\label{sf19:remaining-support}
前系の仮定の下で，非空の高分母 support の必要候補は次の六型に限られる．
標識を保つ候補は百二十六個である．$\overline H=K_5\setminus H$ とし，
補 graph の表示では孤立頂点を省略する．
\begin{center}
\small
\begin{tabular}{crcrcr}
型&$|E|$&降順の次数列&$b_1$&$\overline H$&標識付き個数\\\hline
$G_{5,11}$&7&$(4,3,3,2,2)$&3&$P_4$&60\\
$G_{5,13}$&7&$(4,4,2,2,2)$&3&$K_3$&10\\
$G_{5,14}$&8&$(4,3,3,3,3)$&4&$2K_2$&15\\
$G_{5,15}$&8&$(4,4,3,3,2)$&4&$P_3$&30\\
$G_{5,16}$&9&$(4,4,4,3,3)$&5&$K_2$&10\\
$G_{5,17}$&10&$(4,4,4,4,4)$&6&$\varnothing$&1
\end{tabular}
\end{center}
識別子は旧版から固定し，欠番を詰めない．
これらは必要条件を通過した graph の個数であり，
整数接着データ，作用 atlas，または Lagrangian fibrations の存在数を表さない．
\end{corollary}

\begin{proof}
全 $2^{10}=1024$ 個の標識付き部分 graph に既述の障害を適用する．
今回排除された $G_{5,7},G_{5,9}$ は自己同型群の位数がともに二であり，
それぞれ $5!/2=60$ 標識を持つ．従って残りは
\[
 246-60-60=126
\]
である．六代表の自己同型群の位数は，表の順に $2,12,8,4,12,120$ となり，
軌道数 $5!/|\operatorname{Aut}(H)|$ は表の個数と一致する．
各代表の補 graph を取ると第五列が得られる．
\nolinkurl{verify_support_agenda_v19.py} は全 graph に対して，
旧識別子と軌道数，二価頂点の三角形条件，各頂点の削除後の連結性を再確認し，
\nolinkurl{support_agenda_verification_v19.json} に全標識を保存する．
この graph 列挙は先行する整数接着の排除定理を適用するものであり，
それらの証明の代わりではない．
\end{proof}

残る低辺は高々三本である．七高辺の二候補は，低辺が道 $P_4$ を作る型と，
低三角形を作る book 型に分かれる．後者の低三角形には
定理~\ref{bk19:finite-completion} の固定 core に対する有限完成法を適用できるが，
固定 core 自体の無限集合は未処理である．
全二価頂点の高近傍は互いに接続しているため，負の二高分母枝を先に捨ててはならない．
三高分母 cores には完全 matching 定理を用い，四高分母 cores へは無条件で拡張しない．

$G_{5,15},G_{5,16},G_{5,17}$ は $K_4$ を部分 graph として含む．
四使用頂点 $K_4$ や $K_4+\mathrm{leaf}$ の排除は support 全体に関する定理なので，
これらの密な graph を同じ理由で除くことはできない．
一般の五次分母剛性と，六候補の整数接着・幾何的実現は未解決である．
